\section{Phụ lục về AI tạo sinh và cách kiểm chứng}

\subsection{Lập trình, tích hợp và soạn báo cáo}

Trong quá trình thực hiện, nhóm tự thiết kế kiến trúc ranh giới công khai/ẩn, đặc tả luật và trực tiếp cài đặt các thuật toán lõi (phân giải vùng, đánh giá ngữ nghĩa, mã hóa CNF và bộ giải DPLL). Nhóm sử dụng Codex và Antigravity như công cụ hỗ trợ gợi ý khung kiểm thử, đối chiếu biên điều kiện và rà soát tính nhất quán cho giao diện cũng như báo cáo. Khi đặt yêu cầu cho công cụ, nhóm luôn nêu rõ ranh giới dữ liệu công khai, đầu ra mong đợi và các trường hợp phải từ chối.

Chẳng hạn, khi rà soát phần suy luận, nhóm đặt yêu cầu: “Chỉ truyền \texttt{PublicState} cho tác tử; với mỗi phán quyết, thử giả định phủ định và làm rõ điều kiện trả về \texttt{UNKNOWN} hoặc \texttt{INCONSISTENT}.” Tương tự, với nhóm luật mở rộng, yêu cầu nêu đủ cả sáu toán tử so sánh và cho phép \texttt{CONDITIONAL} nhận hai cây manh mối bất kỳ. Mọi phần mã và gợi ý thu được đều do nhóm kiểm chứng lại bằng phản ví dụ, đối chiếu phép gán hoàn chỉnh, bộ kiểm thử tự động và nhật ký thực nghiệm chạy trực tiếp trên máy.

\subsection{Đồ họa, sửa ảnh có mặt nạ và nghiệm thu}

Nhật ký \texttt{docs/AI-PROMPTS.md} lưu các prompt tạo atlas, ảnh tham chiếu, lỗi cần chỉnh sửa và thứ tự khung hình. Ở bản phác thảo biểu cảm ban đầu (9 khung cho Ada), nhóm yêu cầu khóa căn chỉnh khung hình với đoạn trích nguyên văn: “keep face center, eye line, head scale, hair silhouette and shoulders in exactly same position across all nine frames”, nhằm tránh tình trạng nhân vật bị nhảy vị trí khi chuyển biểu cảm.

Khi chuyển sang phong cách pixel chính thức với 7 khung hình cho mỗi nhân vật, nhóm chỉ định rõ số hàng/cột, nền RGBA trong suốt, phần trang phục thay đổi, biên mỗi ô và vùng cần để trống cho ImageGen. Với hiệu ứng chớp mắt, Stable Diffusion 1.5 Inpainting chỉ tái tạo vùng mặt nạ quanh mắt, giữ nguyên pixel ngoài mặt nạ và kênh alpha gốc \cite{inpainting}. Khi cập nhật lại 7 khung của nhân vật Finn, nhóm đối chiếu 24 hàng nhân vật còn lại theo từng byte, đồng thời kiểm tra mép alpha và chiều cao tư thế. Toàn bộ ảnh giao diện ở mục 6 là ảnh chụp trực tiếp từ chương trình đang chạy.

\subsection{Thuyết minh và nguồn đối chiếu}

ElevenLabs TTS chuyển kịch bản tiếng Việt thành từng đoạn thoại để khớp với thao tác trên bàn. Quy trình đo thời lượng từng tệp âm thanh, ghép đoạn vào trục hình 24 khung/giây và dùng các khoảng ngắt giọng để căn phụ đề. Các cảnh chọn màn, gợi ý, phán quyết bị từ chối và tự giải được kiểm tra theo trạng thái cửa sổ và thông báo của engine. Âm thanh và hình minh họa phục vụ trình bày, còn phán quyết trong phần mềm luôn xuất phát từ DPLL.

Tài liệu \texttt{docs/REFERENCES.md} ghi phạm vi đã khảo sát. Nhóm tham khảo hướng dẫn thiết kế giao diện trò chơi \cite{gamedev}, đọc phần mô hình và điều phối của clues-by-sam-solver \cite{safina}, đồng thời khảo sát README và cấu trúc hai dự án Griductive \cite{trantai,wadew}. Khi đối chiếu một ý tưởng từ nguồn ngoài, nhóm dùng luật của đề và kết quả thử trên chính mã Python của nhóm để quyết định cách áp dụng.
